\bookchapter{Глава 4: Установка и управление приложениями — «Набор навыков» телефона}

\chapterintro{Приложения делают телефон всемогущим, но ими нужно правильно управлять. Как вы выделяете рабочее место и распределяете права новому сотруднику, так и установка приложений требует соблюдения регламента.}

\booksection{Официальные магазины: самый безопасный «супермаркет приложений»}

\begin{itemize}
  \item \textbf{Android}: Google Play (за рубежом) или фирменные магазины производителей (например, Huawei AppGallery, Xiaomi GetApps).
  \item \textbf{iOS}: только App Store.
  \item \textbf{HarmonyOS}: Huawei AppGallery.
\end{itemize}

Ключевое значение для безопасности: прошедшие проверку, с гарантированными обновлениями, безопасные платежи. Руководство к действию:\textbf{приоритетно используйте только официальные магазины для загрузки!}

\booksection{Режим чистоты / защиты: включаем «антивирусную стену» телефона}

Основные китайские бренды (Huawei, Xiaomi, OPPO, vivo, Honor) имеют встроенный\textbf{Режим чистоты}или\textbf{Режим защиты}— это системный переключатель безопасности.\textbf{Настоятельно рекомендуем включить его всем пользователям}.

\booksubsection{Что это вообще такое?}

Режим чистоты (также называемый «Режим защиты»)\textbf{блокирует установку сторонних приложений}: когда он включён, система\textbf{разрешает установку приложений только из официального магазина}, а установочные пакеты, загруженные из браузера, по QR-кодам или из мессенджеров, будут напрямую заблокированы, пресекая вредоносное ПО в зародыше.

\booksubsection{Для кого подходит}

\begin{itemize}
  \item \textbf{Пожилые пользователи}: не различают «официальную» и «поддельную» рекламу — режим чистоты становится самым надёжным оберегом.
  \item \textbf{Пользователи, нечувствительные к разрешениям приложений}: не нуждаются в установке нишевых приложений, бета-версий, игр для зарубежных серверов.
  \item \textbf{Все пользователи, желающие повысить безопасность устройства}: даже энтузиастам лучше держать его включённым — если возникнет временная необходимость, можно отдельно отключить.
\end{itemize}

\booksubsection{Путь включения (на примере Huawei/Xiaomi)}

\begin{itemize}
  \item \textbf{Huawei/HarmonyOS}: [Настройки] > [Система и обновления] >\textbf{[Режим чистоты]}, выберите «Включить».
  \item \textbf{Xiaomi}: [Настройки] > [Защита конфиденциальности] >\textbf{[Защита конфиденциальности]}или [Настройки] > [Настройки приложений] > [Установка приложений], отключите «Установка неизвестных приложений».
  \item \textbf{OPPO/vivo}: [Настройки] > [Безопасность] >\textbf{[Установка приложений из внешних источников]}, поочерёдно отключите разрешения.
\end{itemize}

\booksubsection{Механизмы безопасности магазинов приложений}

Магазин приложений — не «декорация», за ним стоит полноценный конвейер безопасности:

\begin{itemize}
  \item \textbf{Антивирусное сканирование}: каждое приложение перед публикацией проходит совместную проверку несколькими антивирусными движками для выявления известных троянов и вирусных сигнатур.
  \item \textbf{Проверка разрешений}: магазин проверяет, «обоснованы» ли запрашиваемые приложением разрешения — приложение-фонарик требует доступ к контактам? Прямой отказ или удаление из магазина.
  \item \textbf{Верификация разработчика}: для публикации требуется квалификация предприятия, удостоверение личности или верификация банковской карты. При проблеме можно найти ответственного — это не «серое» приложение.
  \item \textbf{Постоянный мониторинг}: уже опубликованные приложения проходят повторное сканирование, вновь обнаруженное вредоносное поведение ведёт к принудительному удалению.
\end{itemize}

\booksection{Сторонние источники и «электронные паразиты»: крайне высокий риск!}

Загрузка приложений с веб-страниц или по неизвестным ссылкам может содержать вирусы, похищать личные данные, воровать платёжные средства. Здесь мы особо вводим новое понятие:\textbf{«электронные паразиты»}.

\textbf{Что такое «электронные паразиты»?}\par

Это\textbf{«взломанные» версии программ со встроенными троянами}— вы думаете, что «сэкономили», а на самом деле «попались». Распространённые формы:

\begin{itemize}
  \item «Взломанные», «бесплатные», «без рекламы» версии платных программ;
  \item Серые инструменты вроде «накрутки лайков», «захвата красных конвертов», «игровых читов»;
  \item Так называемые «внутренние версии» приложений из рекламы в групповых чатах.
\end{itemize}

Эти программы внешне работают, но в фоне могут\textbf{воровать контакты, читать SMS, прослушивать микрофон, майнить криптовалюту}. Режим чистоты —\textbf{первая линия обороны}против них: просто не даёт установить.

\booksubsection{Рекомендации второй линии обороны:}

\begin{itemize}
  \item \textbf{Регулярное полное антивирусное сканирование}: используйте встроенный центр безопасности телефона (например, Huawei Phone Manager, Xiaomi Security Center) для ежемесячной проверки.
  \item \textbf{Осторожно с чувствительными разрешениями}: контакты, SMS, местоположение, камера, микрофон — эти пять «высокорисковых разрешений» требуют строгой проверки.
  \item \textbf{Следите за аномалиями расхода батареи и трафика}: телефон внезапно греется, быстро разряжается, трафик уходит по ночам — скорее всего, в фоне прячется «паразит».
\end{itemize}

Если без стороннего приложения никак не обойтись:

\begin{enumerate}
  \item Убедитесь в абсолютной надёжности источника.
  \item Перед установкой внимательно проверьте запрашиваемые разрешения.
  \item После установки просканируйте приложение антивирусом.
\end{enumerate}

Главный совет: отключите в системных настройках разрешение «Установка неизвестных приложений».

\booksection{Установка, обновление, удаление}

\begin{itemize}
  \item \textbf{Установка}: в магазине нажмите «Получить» или «Установить», после установки при первом запуске приложение запросит разрешения.
  \item \textbf{Обновление}:\textbf{обязательно обновляйте!}Главное — исправление уязвимостей безопасности. Рекомендуется включить автоматическое обновление.
  \item \textbf{Удаление}: долгое нажатие на иконку и перетаскивание в зону удаления. Для более полной очистки зайдите в [Настройки] > [Управление приложениями], найдите и удалите.
\end{itemize}

\booksection{Управление разрешениями: вы — «начальник» приложений}

Приложения запрашивают разрешения (например, контакты, местоположение). Иногда это обоснованно (картам нужно местоположение), иногда нет (фонарику не нужны контакты).

\begin{itemize}
  \item \textbf{Варианты разрешений}: «Всегда», «Только при использовании», «Запретить», «Спрашивать каждый раз».
  \item \textbf{Принципы управления}:
  \begin{enumerate}
    \item Минимизация: не давайте без необходимости.
    \item Последующий аудит: регулярно проверяйте в настройках разрешений, кто «превышает полномочия».
    \item Контекстность: например, приложению такси давайте местоположение только при использовании.
  \end{enumerate}
\end{itemize}

\begin{tcbwarn}{Разрешения под пристальным вниманием}
\begin{itemize}
  \item Контакты / SMS: могут использоваться для рассылки спама.
  \item Местоположение в реальном времени: может отслеживаться ваше передвижение.
  \item Камера / Микрофон: возможна скрытая съёмка и запись.
  \item Доступ к файлам: может читать ваши личные файлы.
\end{itemize}

\end{tcbwarn}

\booksection{Состав установочного пакета приложения}

Вы когда-нибудь задумывались, что находится внутри установочного пакета, скачанного из магазина приложений? Понимание состава установочного пакета поможет лучше понять, как работают приложения, и знать, с чего начать поиск при возникновении проблем.

\booksubsection{Структура APK-файла Android}

APK (Android Package Kit) по сути является ZIP-архивом. После распаковки можно увидеть следующие ключевые каталоги и файлы:

\begin{itemize}
  \item \textbf{AndroidManifest.xml}: «удостоверение личности» и «инструкция» приложения. Содержит имя пакета, версию, необходимые разрешения, поддерживаемые версии системы, точку входа Activity и другую ключевую информацию. Система читает этот файл первой при установке.
  \item \textbf{classes.dex}: основной код приложения. Исходный код на Java/Kotlin после компиляции превращается в байт-код Dalvik и хранится в этом файле. Приложение может содержать несколько dex-файлов (например, classes2.dex, classes3.dex).
  \item \textbf{resources.arsc}: скомпилированная индексная таблица ресурсов. Содержит идентификаторы всех ресурсов и их соответствия.
  \item \textbf{Каталог res/}: папка ресурсов. Содержит изображения, файлы макетов, строки, определения цветов и т.д. Изображения для разных разрешений хранятся в разных папках drawable (например, drawable-hdpi, drawable-xhdpi).
  \item \textbf{Каталог assets/}: папка исходных ресурсов. Хранит ресурсы, не проходящие компиляцию, такие как веб-файлы, шрифты, аудио. Приложение может читать их через AssetManager во время выполнения.
  \item \textbf{Каталог lib/}: папка нативных библиотек. Хранит скомпилированные .so-файлы (результат компиляции C/C++ кода), отсортированные по архитектуре CPU (например, arm64-v8a, armeabi-v7a). Игры и приложения, требующие высокой производительности, обычно содержат этот каталог.
  \item \textbf{Каталог META-INF/}: папка с информацией о подписи. Содержит файлы MANIFEST.MF, CERT.SF, CERT.RSA для проверки целостности и подписи APK.
\end{itemize}

\booksubsection{Структура IPA-файла iOS}

IPA (iOS App Store Package) также является архивом, но структура несколько отличается:

\begin{itemize}
  \item \textbf{Каталог Payload/}: основное содержимое. Внутри находится папка .app, содержащая всё содержимое приложения.
  \item \textbf{Info.plist}: конфигурационный файл приложения, аналог AndroidManifest.xml в Android. Содержит ID приложения, версию, требования к разрешениям, поддерживаемые устройства и т.д.
  \item \textbf{Исполняемый файл}: основная программа с именем, совпадающим с именем папки .app, представляет собой скомпилированный бинарный файл в формате Mach-O.
  \item \textbf{Каталог Frameworks/}: динамические библиотеки и фреймворки, от которых зависит приложение.
  \item \textbf{Файлы ресурсов}: изображения, аудио, видео и другие ресурсы, обычно хранятся внутри папки .app.
  \item \textbf{Каталог \_CodeSignature/}: информация о подписи кода, используется для проверки целостности и источника приложения.
\end{itemize}

\booksubsection{Механизм проверки подписи}

Почему приложения из официальных магазинов безопаснее? Проверка подписи — ключевой фактор.

\begin{itemize}
  \item \textbf{Подпись Android}: разработчик подписывает APK своим закрытым ключом, система при установке проверяет открытым ключом. Если APK был изменён (например, в него внедрён вирус), подпись становится недействительной, и система откажет в установке. Google Play также требует App Signing, при котором Google хранит ключ подписи, дополнительно защищая безопасность приложения.
  \item \textbf{Подпись iOS}: проверка Apple ещё строже. Каждый IPA должен быть подписан сертификатом, выданным Apple, иначе установка невозможна. Именно поэтому iOS позволяет установку только из App Store или через корпоративный сертификат. Срок действия приложений, подписанных корпоративным сертификатом, ограничен, и Apple может отозвать его в любой момент.
  \item \textbf{Функции подписи}:
  \begin{enumerate}
    \item Проверка источника: подтверждение, что приложение действительно от заявленного разработчика.
    \item Проверка целостности: гарантия, что приложение не было изменено при передаче.
    \item Защита от подделки: вредоносное приложение не может выдать себя за другое.
  \end{enumerate}
\end{itemize}

\booksubsection{Размер установочного пакета и оптимизация}

Вы могли заметить, что размер приложения, показанный в магазине, отличается от занимаемого после установки. Это потому, что:

\begin{itemize}
  \item \textbf{Размер загрузки}: размер сжатого APK/IPA.
  \item \textbf{Размер установки}: размер после распаковки, обычно значительно больше размера загрузки.
  \item \textbf{Занимаемое при работе}: во время работы приложение создаёт кэш и файлы данных, занимая дополнительное пространство.
\end{itemize}

Современные методы оптимизации приложений:

\begin{itemize}
  \item \textbf{App Bundle (Android)}: технология динамического распространения Google Play. При загрузке пользователь получает только ресурсы, подходящие для его устройства (например, изображения только нужного разрешения), что значительно уменьшает размер загрузки.
  \item \textbf{App Thinning (iOS)}: аналогичная технология, Apple автоматически обрезает ненужные ресурсы в зависимости от устройства.
  \item \textbf{Загрузка по требованию}: крупные приложения, такие как игры, могут сначала загрузить основной пакет, а затем докачать пакеты ресурсов.
\end{itemize}

\begin{tcbinfo}{Маленькая заметка: почему некоторые установочные пакеты такие большие?}
Игровые приложения обычно самые большие, так как содержат множество текстур высокого разрешения, 3D-моделей и аудиофайлов. Социальные приложения тоже немаленькие, потому что поддерживают множество функциональных модулей. Утилиты обычно самые маленькие. Если вы заметили, что какая-то утилита подозрительно большая — стоит насторожиться, не упакованы ли в неё ненужные функции.

\end{tcbinfo}

\begin{tcbpractice}{Практическое задание: проверьте разрешения 3 приложений}
Выберите 3 часто используемых приложения (например, соцсеть, магазин, утилиту), зайдите в [Настройки] на страницу управления разрешениями приложений, проверьте каждое и отключите ненужные чувствительные разрешения или переключите на «Только при использовании».

\end{tcbpractice}
